% =========================================================================
%  Electromagnetismo I — Tema 0: Herramientas matemáticas
%  Estilo: artículo científico minimalista, heredado del informe «Efecto
%  Doppler con ultrasonidos» (amsart, Palatino, versalitas, enunciados
%  numerados por sección, paleta de azules y grises sobrios), con la
%  notación vectorial de «Cálculo integral multivariable» (L. Navas) y las
%  expresiones del formulario del curso.
% =========================================================================

% --- IDIOMA Y CODIFICACIÓN ------------------------------------------------
\usepackage[utf8]{inputenc}
\usepackage[T1]{fontenc}
\usepackage[spanish,es-tabla,es-nodecimaldot]{babel}

% --- TIPOGRAFÍA (Palatino, como el documento de referencia) ---------------
\usepackage{mathpazo}
\linespread{1.05}
\usepackage[final]{microtype}

% --- MATEMÁTICAS -------------------------------------------------------------
\usepackage{amsmath, amssymb, mathtools}
\usepackage{bm}
\usepackage{mathrsfs}

% --- PÁGINA, COLORES Y GRÁFICOS --------------------------------------------
\usepackage[a4paper,margin=2.6cm]{geometry}
\usepackage{graphicx}
\usepackage{xcolor}
% Paleta base (idéntica a la del informe de ecuaciones diferenciales)
\definecolor{azulprof}{HTML}{1B365D}    % azul marino
\definecolor{azulacero}{HTML}{3B6EA8}   % azul acero
\definecolor{azulclaro}{HTML}{B7CBE3}   % azul hielo
\definecolor{grispizarra}{HTML}{5B6770} % gris pizarra
\definecolor{grisclaro}{HTML}{E4E8EC}   % gris perla
% Paleta ampliada para campos escalares y vectoriales (magnitud creciente:
% índigo → púrpura → carmín → naranja → ámbar), sobria y legible en impresión
\definecolor{cm0}{RGB}{38,30,112}
\definecolor{cm1}{RGB}{112,36,140}
\definecolor{cm2}{RGB}{190,52,118}
\definecolor{cm3}{RGB}{236,104,66}
\definecolor{cm4}{RGB}{252,184,42}
% Colores de apoyo para diagramas geométricos
\definecolor{rojoteo}{HTML}{B3262E}     % trayectorias y circuitos
\definecolor{verdeteo}{HTML}{2E8B57}    % vector unitario / normal
\definecolor{ocreteo}{HTML}{C98A12}     % superficies y volúmenes

\usepackage{tikz}
\usepackage{tikz-3dplot}
\usetikzlibrary{babel, arrows.meta, calc, positioning, decorations.pathmorphing,
                decorations.markings, patterns, 3d, intersections, angles, quotes,
                fadings, shapes.geometric}
\usepackage{pgfplots}
\pgfplotsset{compat=1.17}
\usepgfplotslibrary{colormaps, patchplots, fillbetween}
\pgfplotsset{colormap={campo}{
  rgb255=(38,30,112) rgb255=(112,36,140) rgb255=(190,52,118)
  rgb255=(236,104,66) rgb255=(252,184,42) rgb255=(255,226,110)}}
\pgfplotsset{colormap={campoclaro}{
  rgb255=(232,236,240) rgb255=(183,203,227) rgb255=(59,110,168) rgb255=(27,54,93)}}
\pgfplotsset{colormap={verdeamarillo}{
  rgb255=(120,200,70) rgb255=(190,220,80) rgb255=(250,225,100) rgb255=(250,190,90)}}
\usepackage{booktabs}
\usepackage{array}
\usepackage{enumitem}
\setlist[itemize]{leftmargin=1.6em, itemsep=2pt, topsep=4pt}
\setlist[enumerate]{leftmargin=1.8em, itemsep=2pt, topsep=4pt}

% --- ENLACES Y REFERENCIAS CRUZADAS ----------------------------------------
\usepackage[colorlinks=true, linkcolor=azulacero, citecolor=azulacero,
            urlcolor=azulacero, plainpages=false, pdfusetitle=false,
            pdftitle={Tema 0. Herramientas matemáticas},
            pdfauthor={Daniel Martín Guerrero}]{hyperref}
\usepackage[spanish,nameinlink]{cleveref}

% --- BIBLIOGRAFÍA -----------------------------------------------------------
\usepackage{csquotes}
\usepackage[backend=biber, style=numeric, sorting=none, giveninits=true,
            maxbibnames=6, minbibnames=3, doi=true, url=true,
            isbn=false]{biblatex}
\addbibresource{referencias.bib}
\DeclareFieldFormat[article,techreport,unpublished,book]{title}{#1\isdot}
\DeclareFieldFormat{url}{\url{#1}}
\renewcommand*{\bibfont}{\small}
\setlength{\bibitemsep}{3pt}

% --- CABECERAS ---------------------------------------------------------------
\pagestyle{headings}
\makeatletter
\def\ps@headings{%
  \let\@mkboth\markboth
  \def\@evenhead{\small\scshape\normalfont\thepage\hfil Electromagnetismo I}%
  \def\@oddhead{\small\scshape\normalfont Tema 0. Herramientas matemáticas\hfil\thepage}%
  \let\@evenfoot\@empty \let\@oddfoot\@empty}
\makeatother
\numberwithin{equation}{section}
\setcounter{tocdepth}{2}
\makeatletter
\def\@@and{y}
\renewcommand{\keywordsname}{Palabras clave}
\makeatother

% =========================================================================
%               NOTACIÓN Y MACROS DE USO FRECUENTE
% =========================================================================
% --- Diferencial, derivadas parciales y números -----------------------------
\newcommand{\dif}{\mathrm{d}}
\newcommand{\pd}[2]{\frac{\partial #1}{\partial #2}}
\newcommand{\pdd}[2]{\frac{\partial^{2} #1}{\partial #2^{2}}}
\newcommand{\pdm}[3]{\frac{\partial^{2} #1}{\partial #2\,\partial #3}}
\newcommand{\R}{\mathbb{R}}
\newcommand{\N}{\mathbb{N}}
\newcommand{\Z}{\mathbb{Z}}
\newcommand{\eul}{\mathrm{e}}

% --- Vectores, versores y operador nabla ---------------------------------
\newcommand{\vc}[1]{\vec{#1}}                       % vector con flecha
\newcommand{\uh}[1]{\hat{u}_{#1}}                   % versor û_#1
\newcommand{\ux}{\uh{x}}\newcommand{\uy}{\uh{y}}\newcommand{\uz}{\uh{z}}
\newcommand{\ur}{\uh{r}}\newcommand{\uth}{\uh{\theta}}\newcommand{\uph}{\uh{\varphi}}
\newcommand{\urho}{\uh{\rho}}\newcommand{\uR}{\uh{R}}\newcommand{\un}{\hat{n}}
\newcommand{\nab}{\vec{\nabla}}                     % ∇⃗
\newcommand{\nabp}{\vec{\nabla}'}                   % ∇⃗′ (coordenadas fuente)
\newcommand{\Lap}{\nabla^{2}}                       % laplaciano
\newcommand{\Lapp}{\nabla'^{2}}                      % laplaciano en coordenadas fuente
\newcommand{\norm}[1]{\lVert #1\rVert}
\newcommand{\abs}[1]{\lvert #1\rvert}
\newcommand{\prodesc}[2]{\langle #1,#2\rangle}
\newcommand{\rp}{\vc{r}\,'}                         % r⃗′ punto fuente
\newcommand{\dl}{\dif\vc{\ell}}                     % elemento de línea dirigido
\newcommand{\dS}{\dif\vc{S}}                        % elemento de área dirigido
\newcommand{\dsc}{\dif S}                           % elemento de área escalar
\newcommand{\dV}{\dif V}                            % elemento de volumen
\newcommand{\defeq}{\mathrel{\overset{\mathrm{def}}{=}}}
\newcommand{\bordeS}{\partial S}
\newcommand{\bordeV}{\partial V}
\DeclareMathOperator{\grad}{grad}
\DeclareMathOperator{\dive}{div}
\DeclareMathOperator{\rot}{rot}
\DeclareMathOperator{\sen}{sen}
\DeclareMathOperator{\senh}{senh}
\DeclareMathOperator{\sgn}{sgn}
\DeclareMathOperator{\Tr}{Tr}
\renewcommand{\sin}{\sen}                            % «sen» en todo el texto

% Integral de superficie cerrada (doble integral con círculo)
\newcommand{\oiint}{\mathop{\vcenter{\hbox{$\displaystyle\oint\mkern-13mu\oint$}}}\nolimits}
\newcommand{\oiintt}{\mathop{\vcenter{\hbox{$\textstyle\oint\mkern-11mu\oint$}}}\nolimits}

% Tensor de Levi-Civita y delta de Kronecker
\newcommand{\eps}{\varepsilon}
\newcommand{\dpart}[1]{\partial_{#1}}

% Texto de relleno discreto para datos pendientes
\newcommand{\pend}[1]{\textcolor{grispizarra}{\colorbox{grisclaro}{\strut\,#1\,}}}

% =========================================================================
%        ENTORNOS FORMALES (enunciados numerados, estilo artículo)
% =========================================================================
\usepackage{amsthm}

% Enunciados en cursiva (resultados generales)
\newtheoremstyle{cienplano}{10pt}{10pt}{\itshape}{}%
  {\bfseries\color{azulprof}}{.}{0.6em}%
  {\thmname{#1}\thmnumber{ #2}\thmnote{\normalfont\upshape\ (#3)}}
% Definiciones, cálculos y ejemplos (cuerpo en redonda)
\newtheoremstyle{ciendef}{10pt}{10pt}{\normalfont}{}%
  {\bfseries\color{azulprof}}{.}{0.6em}%
  {\thmname{#1}\thmnumber{ #2}\thmnote{\normalfont\upshape\ (#3)}}
% Observaciones y notas (encabezado en cursiva)
\newtheoremstyle{cienobs}{8pt}{8pt}{\normalfont}{}%
  {\itshape\color{azulprof}}{.}{0.6em}%
  {\thmname{#1}\thmnumber{ #2}\thmnote{\normalfont\upshape\ (#3)}}

\theoremstyle{cienplano}
\newtheorem{proposicion}{Proposición}[section]
\newtheorem{teorema}[proposicion]{Teorema}
\newtheorem{lema}[proposicion]{Lema}
\newtheorem{corolario}[proposicion]{Corolario}
\theoremstyle{ciendef}
\newtheorem{definicion}[proposicion]{Definición}
\newtheorem{notacion}[proposicion]{Notación}
\newtheorem{ejemplo}[proposicion]{Ejemplo}
\newtheorem{ejercicio}[proposicion]{Ejercicio}
\theoremstyle{cienobs}
\newtheorem{observacion}[proposicion]{Observación}

% Resolución de ejercicios: misma maquinaria que «proof», con otro rótulo
\renewcommand{\qedsymbol}{\ensuremath{\blacksquare}}
\newenvironment{resolucion}[1][Resolución]
  {\begin{proof}[\itshape\color{azulprof}#1]}{\end{proof}}

\crefname{proposicion}{proposición}{proposiciones}
\Crefname{proposicion}{Proposición}{Proposiciones}
\crefname{teorema}{teorema}{teoremas}
\Crefname{teorema}{Teorema}{Teoremas}
\crefname{lema}{lema}{lemas}
\Crefname{lema}{Lema}{Lemas}
\crefname{corolario}{corolario}{corolarios}
\Crefname{corolario}{Corolario}{Corolarios}
\crefname{definicion}{definición}{definiciones}
\Crefname{definicion}{Definición}{Definiciones}
\crefname{notacion}{notación}{notaciones}
\Crefname{notacion}{Notación}{Notaciones}
\crefname{ejemplo}{ejemplo}{ejemplos}
\Crefname{ejemplo}{Ejemplo}{Ejemplos}
\crefname{ejercicio}{ejercicio}{ejercicios}
\Crefname{ejercicio}{Ejercicio}{Ejercicios}
\crefname{observacion}{observación}{observaciones}
\Crefname{observacion}{Observación}{Observaciones}
\crefname{section}{sección}{secciones}
\Crefname{section}{Sección}{Secciones}
\crefname{subsection}{sección}{secciones}
\Crefname{subsection}{Sección}{Secciones}
\crefname{table}{tabla}{tablas}
\Crefname{table}{Tabla}{Tablas}
\crefname{figure}{figura}{figuras}
\Crefname{figure}{Figura}{Figuras}
\crefname{equation}{ecuación}{ecuaciones}
\Crefname{equation}{Ecuación}{Ecuaciones}
\crefformat{equation}{ecuación~(#2#1#3)}
\Crefformat{equation}{Ecuación~(#2#1#3)}
\crefrangeformat{equation}{ecuaciones~(#3#1#4)--(#5#2#6)}
\crefmultiformat{equation}{ecuaciones~(#2#1#3)}{ y~(#2#1#3)}{, (#2#1#3)}{ y~(#2#1#3)}

% Pies de figura y tabla, en la línea del documento de referencia
\usepackage[font=small, labelfont={sc,color=azulprof}, labelsep=period,
            margin=0.6cm, justification=justified]{caption}

% =========================================================================
%        ESTILOS GRÁFICOS COMUNES (TikZ / pgfplots)
% =========================================================================
% =========================================================================
%  Estilos y macros gráficas comunes
%  - \mapcolor{t}: escala cromática «campo» (t en [0,1]); define el color mapc
%  - \flechas{u}{v}{lista x}{lista y}{escala}{|F|max}: campo vectorial 2D
%  - \flechasIII{u}{v}{w}{...}: campo vectorial 3D (con tikz-3dplot)
% =========================================================================

% --- Escala cromática (5 tramos entre cm0 y cm4) ------------------------
\newcommand{\mapcolor}[1]{%
  \pgfmathsetmacro{\mt}{min(max(#1,0),0.9999)*4}%
  \pgfmathtruncatemacro{\mi}{floor(\mt)}%
  \pgfmathtruncatemacro{\mj}{\mi+1}%
  \pgfmathsetmacro{\mf}{100*(\mt-\mi)}%
  \colorlet{mapc}{cm\mj!\mf!cm\mi}}

% --- Estilos TikZ generales ------------------------------------------------
\tikzset{
  >=Stealth,
  eje/.style={->, line width=0.6pt, grispizarra},
  ejeprof/.style={->, line width=0.7pt, azulprof},
  vect/.style={->, line width=1pt, azulprof},
  vectrojo/.style={->, line width=1pt, rojoteo},
  vectverde/.style={->, line width=1pt, verdeteo},
  fino/.style={line width=0.4pt, grispizarra},
  punteado/.style={densely dotted, line width=0.5pt, grispizarra},
  etiq/.style={font=\small, text=black!85},
  cuadro/.style={draw=grispizarra!60, line width=0.4pt, fill=white, inner sep=2pt},
  trayecto/.style={line width=1.1pt, rojoteo},
  orienta/.style={decoration={markings, mark=at position 0.55 with {\arrow{Stealth[length=5pt]}}}, postaction={decorate}},
}

% --- Marco con marcas para gráficas de campos 2D ---------------------------
% \marco2D{xmin}{xmax}{ymin}{ymax}{xticks}{yticks}  (unidades TikZ = unidades del campo)
\newcommand{\marcoD}[6]{%
  \draw[grispizarra, line width=0.5pt] (#1,#3) rectangle (#2,#4);
  \foreach \t in {#5}{%
     \draw[grispizarra, line width=0.4pt] (\t,#3) -- ++(0,0.08);
     \draw[grispizarra, line width=0.4pt] (\t,#4) -- ++(0,-0.08);
     \node[below=1pt, font=\scriptsize] at (\t,#3) {\pgfmathprintnumber[fixed,precision=1,use comma,zerofill=false]{\t}};}%
  \foreach \t in {#6}{%
     \draw[grispizarra, line width=0.4pt] (#1,\t) -- ++(0.08,0);
     \draw[grispizarra, line width=0.4pt] (#2,\t) -- ++(-0.08,0);
     \node[left=1pt, font=\scriptsize] at (#1,\t) {\pgfmathprintnumber[fixed,precision=1,use comma,zerofill=false]{\t}};}%
}

% --- Campo vectorial plano: flechas centradas en cada nudo de la malla ---
% #1,#2: componentes u(x,y), v(x,y) (en términos de \x e \y)
% #3,#4: listas de abscisas y ordenadas; #5: factor de escala de longitud
% #6: módulo máximo (normaliza el color)
\newcommand{\flechas}[6]{%
  \foreach \x in {#3}{\foreach \y in {#4}{%
    \pgfmathsetmacro{\fu}{#1}\pgfmathsetmacro{\fv}{#2}%
    \pgfmathsetmacro{\fm}{sqrt(\fu*\fu+\fv*\fv)}%
    \pgfmathsetmacro{\ft}{\fm/(#6)}%
    \mapcolor{\ft}%
    \ifdim\fm pt>0.001pt
      \pgfmathsetmacro{\hl}{2.2+4.2*\ft}%
      \pgfmathsetmacro{\lw}{0.35+0.5*\ft}%
      \draw[-{Stealth[length=\hl pt,width=\hl pt]}, line width=\lw pt, mapc]
        (\x-0.5*#5*\fu,\y-0.5*#5*\fv) -- ++(#5*\fu,#5*\fv);
    \else
      \fill[mapc] (\x,\y) circle (0.7pt);
    \fi}}}

% --- Campo vectorial espacial (tikz-3dplot): lo mismo en 3D ----------------
\newcommand{\flechasIII}[8]{%
  \foreach \x in {#4}{\foreach \y in {#5}{\foreach \z in {#6}{%
    \pgfmathsetmacro{\fu}{#1}\pgfmathsetmacro{\fv}{#2}\pgfmathsetmacro{\fw}{#3}%
    \pgfmathsetmacro{\fm}{sqrt(\fu*\fu+\fv*\fv+\fw*\fw)}%
    \pgfmathsetmacro{\ft}{\fm/(#8)}%
    \mapcolor{\ft}%
    \ifdim\fm pt>0.001pt
      \pgfmathsetmacro{\hl}{2.0+3.6*\ft}%
      \pgfmathsetmacro{\lw}{0.35+0.5*\ft}%
      \draw[-{Stealth[length=\hl pt,width=\hl pt]}, line width=\lw pt, mapc]
        (\x-0.5*#7*\fu,\y-0.5*#7*\fv,\z-0.5*#7*\fw) -- ++(#7*\fu,#7*\fv,#7*\fw);
    \else
      \fill[mapc] (\x,\y,\z) circle (0.7pt);
    \fi}}}}

% --- Estilo base de ejes pgfplots (como en la figura del efecto Doppler) --
\pgfplotsset{
  tema/.style={
    tick label style={font=\scriptsize}, label style={font=\small},
    axis lines=left, axis line style={grispizarra},
    grid=major, grid style={grisclaro, line width=0.4pt},
    every tick/.style={grispizarra},
    legend style={font=\scriptsize, draw=none, fill=white, fill opacity=0.85,
                  text opacity=1},
    /pgf/number format/use comma,
  },
  mapa/.style={
    tick label style={font=\scriptsize}, label style={font=\small},
    axis line style={grispizarra}, every tick/.style={grispizarra},
    /pgf/number format/use comma,
    colormap name=campo,
  },
  superficie/.style={
    tick label style={font=\scriptsize}, label style={font=\small},
    axis line style={grispizarra}, every tick/.style={grispizarra},
    grid=none, /pgf/number format/use comma,
    colormap name=campo,
  },
}

% Rótulos de ejes para marcoD:  \ejesD{xmin}{xmax}{ymin}{ymax}
\newcommand{\ejesD}[4]{%
  \node[font=\small, anchor=north] at ($({(#1+#2)/2},#3)+(0,-13pt)$) {$x$ (u.\,arb.)};
  \node[font=\small, rotate=90, anchor=south] at ($(#1,{(#3+#4)/2})+(-17pt,0)$) {$y$ (u.\,arb.)};}

% Recorta el contenido #5 al rectángulo [#1,#2]x[#3,#4]
\newcommand{\clipD}[5]{\begin{scope}\clip (#1,#3) rectangle (#2,#4); #5 \end{scope}}

% Versor dibujado en el punto (#1,#2,#3) con desplazamiento #4 (coordenadas 3D)
\newcommand{\uvect}[5]{%  versor en P: origen (#1,#2,#3), dirección (#4) y rótulo (#5)
  \draw[-{Stealth[length=4.5pt,width=4pt]}, rojoteo, line width=1.1pt] (#1,#2,#3) -- ++#4 node[#5] {};}

% Mapa de calor de f=exp(-r^2) (anillos concéntricos) con curvas de nivel y marco; #1 = extras
\newcommand{\mapaGauss}[1]{%
  \begin{scope}
    \clip (-2,-2) rectangle (2,2);
    \foreach \k in {0,...,90}{%
      \pgfmathsetmacro{\rk}{3.0*(1-\k/91)}%
      \pgfmathsetmacro{\fk}{exp(-\rk*\rk)}%
      \mapcolor{\fk}%
      \fill[mapc] (0,0) circle (\rk);}
    \foreach \lv in {0.9,0.8,0.7,0.6,0.5,0.4,0.3,0.2,0.1}{%
      \pgfmathsetmacro{\rl}{sqrt(-ln(\lv))}%
      \draw[black!75, line width=0.45pt] (0,0) circle (\rl);}
    #1
  \end{scope}
  \marcoD{-2}{2}{-2}{2}{-2,-1,0,1,2}{-2,-1,0,1,2}
  \ejesD{-2}{2}{-2}{2}}


